%%%%%
%
\section{Moyenne-Echantillonnage 
 en tol�rant des valeurs manquantes (param\`etre {\tt short})}

\itemstyle{Op\'erateur~:} \operfmt{ MoyEch\_mnq\_int}
\index{MoyEch\_mnq\_int}

%%%%%OBJET%%%%
\itemstyle{Objet~:}
Moyenne glissante et \'echantillonnage d'un
param\`etre \`a une fr\'equence donn\'ee, diviseur
de la fr\'equence du param\`etre d'entr\'ee.
Le calcul tient compte de la pr�sence des valeurs manquantes
et permet � l'utilisateur de choisir la tol�rance en terme 
de nombre de valeurs manquantes.
Le r�sultat est un param\`etre r\'eel arrondi \`a la valeur
enti\`ere la plus proche. 

%%%%%ARGUMENTS%%%%
\bigskip
%%%%%ENTREES%%%%
\begin{tabarg}
\multicolumn{3}{l}{\titretab {Entr\'ees}}\\
\hline
$X$ & Vecteur, \texttt{short} & Param\`etre d'entr\'ee\\
$R_f$ & Coeff., \texttt{double} & Rapport des fr\'equences $ F_{sortie}\over
F_{entree}$ , dans l'intervalle $]0, 1]$\\
$Decal$ & Coeff., \texttt{int} & D\'ecalage (\textit{nombre de points}), dans l'intervalle $[-1/R_f, 1/R_f]$\\
$Toler\_mnq$ & Coeff., \texttt{int} & Tolerance pour les valeurs manquantes (\textit{nombre de points}),
 entre $0$ et $1/R_f$\\
\end{tabarg}

\medskip
%%%%%SORTIES%%%%
\begin{tabarg}
\multicolumn{3}{l}{\titretab {Sorties}}\\
\hline
$X_m$ & Vecteur, \texttt{float} & R\'esultat nombre entier (m\^eme unit\'e que $X$)\\
\end{tabarg}


%%%%%FORMULE%%%%
\itemstyle{Formule}

%%Moy
$$
Xm[i] = R_f . \sum_{j = i/R_f + Decal}^{(i+1)/R_f -1 + Decal} X[j]
\qquad i=0,Nbre\_Echant(X). R_f
$$


%%%%COMMENTAIRE%%%%
\itemstyle{Commentaire~:}

$R_f$ doit \^etre l'inverse d'un entier. \newline
Le nombre d'\'echantillons du vecteur de sortie est le produit du
nombre d'\'echantillons du vecteur d'entr\'ee par $R_f$, arrondi
\`a l'entier sup\'erieur. \newline

%%%%%AUTEUR%%%%
\itemstyle{Auteur~:}
CNRM/GMEI/TRAMM

%%%%%REFERENCE%%%%
\itemstyle{R\'ef\'erences~:}
